\section{Related Work}
\label{sec:related_work}

Weight-space model merging has gained intense traction due to its ability to consolidate task-specific capabilities into a single parameter set without the prohibitive cost of multi-task fine-tuning or joint training. In this section, we situate our work within the landscape of model merging paradigms and highlight the emerging trend toward complexity that we aim to deconstruct.

\subsection{Linear and Quadratic Model Merging}
The seminal work on weight-space interpolation showed that fine-tuned weights often reside in a shared low-loss basin, allowing simple linear interpolation of checkpoints \cite{Wortsman2022}. This concept was formalized as \textbf{Task Arithmetic} by \citet{Ilharco2022}, who demonstrated that subtracting the pre-trained base model weights from task-specific checkpoints yields ``task vectors'' that can be linearly added, scaled, or negated to edit model capabilities. 

To refine the linear interpolation, several methods propose non-uniform scaling of parameters. \textbf{Fisher Merging} \cite{Matena2022} uses the diagonal elements of the Fisher Information Matrix as weights for parameter fusion, prioritizing coordinates that are crucial for specific tasks. RegCalMerge \cite{RegCalMerge2024} introduces offline calibration and regularization to address transductive overfitting and sacrificial task bias in layer-wise merging. Similarly, PolyMerge and SplineMerge \cite{PolyMerge2025} parameterize merging coefficients as smooth, low-degree polynomials or splines across network depth to regularize the optimization landscape. While successful, these techniques assume a full-density model workspace, leaving them susceptible to parameter interference when scaling to a large number of tasks.

\subsection{Sparse Model Merging and Interference Resolution}
As the number of tasks increases, direct parameter addition is prone to coordinate-wise interference, where opposite sign updates from different tasks cancel each other out. To resolve this, modern research has focused on sparse parameter merging. 

The most prominent framework is \textbf{TIES-Merging} (Trim, Elect, Sign) \cite{Yadav2023}. TIES-Merging argues that parameter interference arises primarily from sign disagreements across different task vectors. To mitigate this, TIES-Merging introduces a coordinate-wise voting heuristic: for each parameter coordinate, tasks vote on whether the update should be positive or negative. The dominant sign is elected, and any task vector that disagrees with this dominant sign has its corresponding coordinate zeroed out. Finally, a disjoint merging is performed on the remaining parameters.

Following TIES-Merging, \textbf{DARE} (Delta-Agnostic Drop and Rescale) \cite{Yu24} introduces a stochastic alternative. Instead of magnitude-based trimming, DARE applies a random dropout mask to task vectors and rescales the surviving updates to preserve the expected value of the updates. Despite its randomized nature, DARE still relies on TIES's sign consensus heuristic to resolve sign conflicts in the final fusion step.

Other works explore the interaction of pruning and merging. ZipMerge \cite{ZipMerge2025} examines joint weight pruning and merging, finding that decoupled ``Prune-then-Merge'' pipelines are often more robust than coupled, active compression frameworks. OFS-Tune \cite{OFS2025} demonstrates that offline few-shot tuning on a small validation set can surpass active online test-time merging (such as AdaMerging \cite{AdaMerging2023}) while avoiding the computational overhead of test-time optimization.

\subsection{Our Departure: The Minimalist Critique}
Unlike previous sparse model merging works that focus on designing new, increasingly complex heuristics, our work is a deconstructive critique. We apply Occam's razor directly to the sign-consensus and sign-voting steps popularized by TIES-Merging. We argue that the success of sparse model merging has been incorrectly attributed to sign-consensus and voting heuristics. Instead, we show that simply applying uniform layer-wise magnitude pruning (to remove fine-tuning noise) followed by direct linear addition (Task Arithmetic) matches the performance of these complex methods (within the margin of statistical error). By demonstrating the redundancy of sign consensus, we restore simplicity and elegance as the guiding principles of model merging.
